
本章我们证明本书所用的若干技术性结果。所有细节都将详尽给出，因为它们必定会阐明
前几章中某些复杂而又必不可少的段落。\footnote{改编自 [52,~Chapter~LXIII]。}

\section{\trFloer{}算子的诸版本}
\label{sec:versions-floer-operator}

在这一技术性的章节里，我们将在记号上区分如下两者：其一是算子
\[
\scr{F}_{\rho}\colon T_{w_{\rho}}\scr{P}^{1,p}(x,z)\longrightarrow L^{p}(\RR\times S^{1};TW)
\]
（它由
\[
\scr{F}_{\rho}(X)=\scr{F}\bigl(\exp_{w_{\rho}}(X)\bigr)=\Bigl(\frac{\partial}{\partial s}+J\frac{\partial}{\partial t}+\operatorname{grad}H\Bigr)\bigl(\exp_{w_{\rho}}(X)\bigr)
\]
来定义），其二是写在诸\trTrivialization{}中的同一算子，我们将把它记作
\[
F_{\rho}\colon W^{1,p}(\RR\times S^{1};\RR^{2n})\longrightarrow L^{p}(\RR\times S^{1};\RR^{2n})\text{。}
\]
我们还要使用（非线性的）算子
\[
F\colon[\rho_{0},+\infty[\,\times\, W^{1,p}(\RR\times S^{1};\RR^{2n})\longrightarrow L^{p}(\RR\times S^{1};\RR^{2n})
\]
它在\cref{sec:construction-psi-rho}中对 $\rho\ge\rho_{0}$ 由 $F(\rho,Y)=F_{\rho}(Y)$ 来定义。我们提醒一下，$F_{\rho}$ 就是\trTrivialization{}之后的算子 $\scr{F}_{\rho}$。

$\scr{F}_{\rho}$ 的源由沿 $w_{\rho}$ 切于 $W$ 而又属于
$W^{1,p}(\RR\times S^{1};\RR^{m})$ 的那些\trVectorField{}构成，这要经由复合
\[
\RR\times S^{1}\longrightarrow TW\hookrightarrow T_{W}\RR^{m}\xrightarrow{\ \ \operatorname{pr}_{2}\ \ }\RR^{m}\text{。}
\]
于是 $\scr{F}_{\rho}(X)$ 是一个沿 $\exp_{w_{\rho}}(X)$ 切于 $W$ 的\trVectorField{}，经由上面的映射，它属于 $L^{p}(\RR\times S^{1};\RR^{m})$。为了定义 $F_{\rho}$，我们使用\cref{sec:construction-psi-rho}所述（并在\cref{fig:trivializations-w-rho}中草绘）的诸\trTrivialization。要验证它是\trWellDefined{}的，只需应用\cref{thm:w1p-product-inequality}便可看出：
\begin{itemize}
\item 若 $y\in W^{1,p}(\RR\times S^{1};\RR^{2n})$，那么 $\sum y_{i}Z_{i}^{\rho}\in T_{w_{\rho}}\scr{P}^{1,p}(x,z)$。
\item 若 $Y\in L^{p}(\RR\times S^{1};TW)$，那么 $\bigl(g(Y,Z_{j}^{\rho,x})\bigr)_{j}\in L^{p}(\RR\times S^{1};\RR^{2n})$。
\end{itemize}
